\begin{proposition}[A Hopf surface degeneration and its comparison map]\hypertarget{tag-0317}{}\label{push:src:B.4.7}
\StatementRecord{push:src:B.4.7}{proved}{0317}{\SeriesRef{push:src:B.4.6}, \SeriesRef{push:src:B.4.2}, \SeriesRef{push:src:A.4.7}; \citep[Definition 1.1, pp.~85--86]{Fujimoto1990AB}; \citet[Section 2.2]{AbashevaDeev2024AB}}

There is a proper semistable analytic family of compact complex
surfaces $p:\mathcal H\to\Delta$ with smooth total space such that:
\begin{enumerate}
\item its punctured restriction is a holomorphic product with the
scalar primary Hopf surface
\[
 H_-=(\C^2\setminus\{0\})/\langle -e^{-\pi}\id\rangle;
\]
\item its central fibre is $Y=H_1\cup_E H_2$, where both $H_i$
are isomorphic to
\[
 H_+=(\C^2\setminus\{0\})/\langle e^{-2\pi}\id\rangle,
 \qquad E=\C/(\Z+i\Z);
\]
\item $N_{E/H_1}\simeq N_{E/H_2}\simeq\mathcal O_E$, and
\[
 \pi_1(Y)=0,\qquad h^1(Y,\mathcal O_Y)=1,
 \qquad H_*(Y;\Z)=(\Z,0,0,\Z^3,\Z^2);
\]
\item no punctured product isomorphism has a meromorphic extension
whose graph is proper over the two total spaces near the central
fibres.
\end{enumerate}
Consequently $\mathcal H\times\PP^1\to\Delta$ is a
semistable family of compact threefolds with simply connected centre
and $h^1(\mathcal O)=1$, whose normalized components fail the
$\partial\bar\partial$ lemma.  Its smooth fibres have Euler number
zero and fundamental group $\Z$.
\end{proposition}


\begin{proof}
\begin{enumerate}[label=\arabic*.,leftmargin=*]
\item We construct the family by an explicit instance of Fujimoto's
generalized logarithmic transformation
\citep[Definition 1.1, pp.~85--86]{Fujimoto1990AB}.
Let
\[
 b:\mathcal B=\operatorname{Bl}_{(p,0)}(\PP^1\times\Delta)
       \longrightarrow\PP^1\times\Delta.
\]
The central curve is $F_1\cup F_2$, where $F_1$ is the strict
transform of $\PP^1\times\{0\}$ and $F_2$ is the exceptional
projective line.  They meet transversely in one point.  The two
line bundles
\[
 L=b^*\operatorname{pr}_1^*\mathcal O_{\PP^1}(1),
 \qquad A=\mathcal O_{\mathcal B}(F_2)
\]
have degree-restriction matrix
\[
 \begin{array}{c|rr}
       &L&A\\\hline
 F_1&1&1\\
 F_2&0&-1
 \end{array}.
\]
In particular this matrix is unimodular.  Take nonzero holomorphic
transition functions $l_{ab},a_{ab}$ for these two line bundles.
On small intersections choose logarithms and put
\[
 t_{ab}=\left[
       \frac{(1+i)\log l_{ab}-i\log a_{ab}}{2\pi i}
                         \right]\in E.
\]
This is independent of the logarithm choices, since $1+i,-i$
belong to the lattice of $E$.  On triple intersections the sum
is again a lattice element.  The $t_{ab}$ therefore form a
holomorphic cocycle with values in the translation group $E$.
Glue $U_a\times E$ by these translations, obtaining a
holomorphic principal $E$-bundle $\mathcal H\to\mathcal B$.
It is a proper submersion, since it is locally a product with a
compact elliptic curve.  Its total space is smooth and Hausdorff.
Compose with $\mathcal B\to\Delta$.

\item The Chern vector of the bundle, in the lattice $\Z+i\Z$, is
\[
 \kappa=(1+i)c_1(L)-i c_1(A).
\]
Its restrictions to $F_1,F_2$ are respectively $1$ and $i$.
On a nonzero fibre its restriction is $1+i$.  These statements
also follow directly from the displayed transition functions.
Locally at the crossing of $F_1$ and $F_2$ the base map is
$xy=u$, and the principal elliptic bundle is locally trivial.
The composed family is thus locally $(x,y,z)\mapsto xy$,
with $z$ a coordinate on $E$.  It is flat, has smooth total
space and reduced simple normal crossings centre.  It is proper
because both maps in the composition are proper.

\item For a primitive vector
$c\in\Z+i\Z$, choose $d$ so that $(c,d)$ is an oriented
lattice basis.  The homomorphism
\[
 \C^*\longrightarrow E,\qquad
 z\longmapsto\left[\frac{c\log z}{2\pi i}\right]
\]
has kernel generated by
$\alpha_c=\exp(2\pi i d/c)$, with $|\alpha_c|<1$.
The principal $E$-bundle on $\PP^1$ associated to
$\mathcal O(1)^\times$ by this homomorphism is the quotient
of $\mathcal O(1)^\times$ by that cyclic kernel.  Fibrewise
inversion identifies $\mathcal O(1)^\times$ with
$\mathcal O(-1)^\times\simeq\C^2\setminus\{0\}$;
inverting the cyclic generator does not change its group.
The resulting surface is the scalar Hopf surface $H_{\alpha_c}$.
For the three vectors in the construction we may take
\[
 (c,d)=(1,i),\quad(i,-1),\quad(1+i,i),
\]
giving respectively
\[
 \alpha_1=\alpha_i=e^{-2\pi},\qquad
 \alpha_{1+i}=\exp\bigl(2\pi i(1+i)/2\bigr)=-e^{-\pi}.
\]
On $\mathcal B|_{\Delta^*}=\PP^1\times\Delta^*$,
$A$ has the canonical nowhere-zero section of its exceptional
divisor.  Its cocycle is therefore a coboundary there, while
$L$ is the constant $\mathcal O(1)$ bundle.  The cocycle coboundary therefore gives a holomorphic product
isomorphism with $H_-\times\Delta^*$.

\item Each central $H_i\to F_i$ is a principal elliptic bundle.
Its intersection with the other component is a smooth fibre,
so a local trivialization over a base disc gives
$N_{E/H_i}=\mathcal O_E$ exactly.  Thus the product of the
two normal lines is trivial, with the smoothing section supplied
by the coordinate $u=xy$.

\item A scalar Hopf surface is
diffeomorphic to $S^1\times S^3$: divide the radial coordinate
by multiplication by $|\alpha|$ and untwist its constant phase
rotation along the radial interval.  Its fibre inclusion has
fundamental-group map
\[
 \Lambda=\Z^2\longrightarrow\Lambda/\Z c_i.
\]
For $c_1=(1,0)$ and $c_2=(0,1)$, these are the two coordinate
quotients.  The van Kampen pushout kills both generators of
$\Lambda$, and hence $\pi_1(Y)=0$.  In the normalization
Mayer--Vietoris sequence the degree-one map, with the second
inclusion negated, is
\[
 H_1(E;\Z)\longrightarrow H_1(H_1;\Z)\oplus H_1(H_2;\Z),
 \qquad \begin{pmatrix}0&1\\-1&0\end{pmatrix}.
\]
It is unimodular.  The degree-zero map is the primitive column
$(1,-1)^t$, the degree-two target is zero, and $H_3(E)=0$.
The exact sequence now gives
\[
 H_1Y=H_2Y=0,\quad
 0\longrightarrow\Z^2\longrightarrow H_3Y
       \longrightarrow\Z\longrightarrow0,
 \quad H_4Y=\Z^2.
\]
All groups displayed in the proposition are therefore free.

\item For coherent cohomology, translations act trivially on
$H^1(E,\mathcal O_E)$, so for either central principal bundle
the sheaf $R^1\pi_*\mathcal O$ on $\PP^1$ is the constant
line bundle.  Its Leray sequence shows that restriction
\[
 H^1(H_i,\mathcal O_{H_i})\longrightarrow H^1(E,\mathcal O_E)
\]
is an isomorphism.  The normalization exact sequence therefore
has degree-one map $(s,t)\mapsto s-t$ from $\C^2$ to $\C$.
The preceding degree-zero map is surjective.  Its kernel has
dimension one, proving $h^1(Y,\mathcal O_Y)=1$.
This also displays exactly why the overlap-injectivity hypothesis
in \ref{push:src:B.4.6} fails: $H^1(E;\C)\to H^1(E,\mathcal O_E)$ has a
one-dimensional kernel.

\item The two scalar Hopf surfaces used above are
not biholomorphic.  An isomorphism lifts to an automorphism of
$\C^2\setminus\{0\}$ that intertwines the two cyclic generators
or their inverses.  Hartogs extension applies to the lift and
its inverse; both fix zero and have invertible derivative there.
Differentiating the intertwining identity gives
$e^{-2\pi}=(-e^{-\pi})^{\pm1}$, which is false.  They have
no rational curves, since a map $\PP^1\to H_\alpha$ lifts to
$\C^2\setminus\{0\}$ and is constant.  The usual resolution
of a bimeromorphic map of smooth compact surfaces therefore
also shows that these two Hopf surfaces are not bimeromorphic.
This agrees with the explicit scalar case of
\citet[arXiv version~1, Section~2.2, pp.~4--5]{AbashevaDeev2024AB}.

\item The fundamental groups also exclude a meromorphic
extension of the punctured product comparison. A sufficiently
small closed semistable neighbourhood in $\mathcal H$ retracts
onto $Y$, by the monomial-gradient argument of \ref{push:src:B.4.2}.  Its
fundamental group is zero.  The product neighbourhood
$H_-\times\Delta$ has fundamental group $\Z$.  If the
punctured comparison had a proper meromorphic graph, resolving
that graph would give a common smooth proper modification of
the two neighbourhoods.  A proper bimeromorphic map between
smooth complex manifolds preserves $\pi_1$, by the
codimension-two avoidance argument in \ref{push:src:A.4.7}.  This is a
contradiction.  Both models already have the indicated
holomorphic map to the same $\PP^1$; the example therefore
retains an extended base map while the total-space comparison
fails.

\item Finally take the product with $\PP^1$.  Its central integral
Betti numbers are
\[
 (1,0,1,3,2,3,2),
\]
and $h^1(\mathcal O)$ remains one.  The nearby fibre is
$H_-\times\PP^1$, of fundamental group $\Z$ and Euler
number zero.  The resulting fundamental group and Euler number are the stated
invariants of these smooth fibres.
\end{enumerate}
\end{proof}